\documentclass{../nndlcourse}

\begin{document}

\title{卷积神经网络 }
\author{数据科学与大数据技术专业}
\date{第5讲 }
\maketitle
\section{卷积神经网络概览}

卷积神经网络 (convolutional neural network, CNN) 是针对图像、语音等具有局部结构数据提出的一类前馈神经网络. 与上一讲的全连接网络相比, CNN 不再把输入简单看成没有内部组织的长向量, 而是把局部连接、权重共享和下采样等先验直接编码到模型结构中.

\begin{enumerate}
    \item \textbf{为什么全连接网络不适合直接处理图像}
    \begin{itemize}
        \item 在全连接层中, 若输入图像大小为 $H \times W$, 通道数为 $C$, 下一层有 $M$ 个神经元, 则仅这一层的参数量就达到
        \begin{equation*}
            HWC \cdot M + M.
        \end{equation*}
        当输入是高分辨率彩色图像时, 这一数量会迅速增大. 例如一幅 $1000 \times 1000$ 的 RGB 图像若直接连接到 $1000$ 个隐藏单元, 参数量约为 $3 \times 10^9$, 训练和存储成本都很高.
        \item 参数量过大不仅带来计算负担, 还会显著提高过拟合风险. 如果样本规模不足, 模型虽然具有很强的表达能力, 却可能只记住训练集中的偶然细节, 无法学到稳定规律.
        \item 更关键的问题在于, 全连接网络会把图像展平成向量, 从而破坏像素之间原有的空间邻接关系. 这样一来, 边缘、纹理、角点等局部模式不再以自然方式出现在模型输入中, 学习过程必须额外花费大量参数去“重新发现”这些结构.
        \item 图像分类往往希望模型对目标的小幅平移、轻微形变具有稳定响应. 这里容易把“平移不变性”说得过强: 单个卷积层更准确地体现的是\textbf{平移等变性}, 即输入平移后, 特征图会发生对应的平移; 真正较强的平移鲁棒性通常还需要结合池化、下采样以及后续分类层共同实现.
    \end{itemize}
    \item \textbf{卷积结构如何降低参数量并保留局部信息}
    \begin{itemize}
        \item CNN 的核心思想是局部连接. 一个卷积神经元只与输入中的一个局部区域相连, 这个局部区域就对应它在当前层看到的感受野. 例如使用 $3 \times 3$ 卷积核时, 单个位置上的输出只依赖其邻域内的 $3 \times 3$ 像素块, 而不是整幅图像.
        \item 对输入通道数为 $C_{\mathrm{in}}$, 输出通道数为 $C_{\mathrm{out}}$ 的二维卷积层, 若卷积核大小为 $K \times K$, 则参数量为
        \begin{equation*}
            K^2 C_{\mathrm{in}} C_{\mathrm{out}} + C_{\mathrm{out}}.
        \end{equation*}
        当 $K = 3$ 时, 每个卷积核在每个输入通道上只含 $9$ 个权重, 这与全连接层中依赖整幅图像大小的参数量形成鲜明对比.
        \item 进一步地, 同一个卷积核会在整张图像上滑动并重复使用, 这就是权重共享. 它意味着“同一种局部模式”无论出现在左上角还是右下角, 都可以由同一组参数来检测. 因此, 参数减少并不是以牺牲结构信息为代价换来的, 反而正好利用了图像中的空间重复性.
        \item 需要注意的是, 深度学习实现中常说的“卷积”多数实际计算的是互相关 (cross-correlation), 即不对卷积核做翻转. 由于卷积核参数本身是可学习的, 这一差异通常不会影响模型表达能力, 因而工程上仍沿用“卷积”这一名称.
    \end{itemize}
    \item \textbf{卷积核为何能够提取边缘等局部特征}
    \begin{itemize}
        \item 卷积核也常称为滤波器 (filter). 其作用可以理解为: 在局部区域内对输入进行加权求和, 从而突出某类变化模式, 同时抑制其他模式. 当滤波器强调强烈变化时, 它往往会对边缘有明显响应; 当滤波器更偏向邻域平均时, 它则起到平滑作用.
        \item 例如对灰度图像使用卷积核
        \begin{equation*}
            \boldsymbol{K} =
            \begin{bmatrix}
                1 & 0 & -1 \\
                1 & 0 & -1 \\
                1 & 0 & -1
            \end{bmatrix},
        \end{equation*}
        若某个局部区域左侧亮、右侧暗, 则加权和会取得较大绝对值, 因而更容易突出竖直边缘. 从这个角度看, 卷积实际上在比较局部区域内部的变化强弱.
        \item 这也解释了为什么早期卷积层往往学习到边缘检测器, 更深层则逐步组合出纹理、部件乃至更抽象的语义特征. 因此, 卷积网络并不是只做“像素压缩”, 而是在逐层构建更有判别力的表示.
        \item 还要注意, 边缘检测核既可以人工设计, 也可以通过反向传播自动学习. 实际深度学习系统更常用后一种方式, 因为不同任务所需的最优局部模式并不完全相同.
    \end{itemize}
    \item \textbf{卷积神经网络的生物学启发与基本结构}
    \begin{itemize}
        \item 卷积神经网络受到生物视觉系统中感受野机制的启发. 在视觉皮层中, 一个神经元通常只对视野中的局部区域敏感; 类似地, CNN 中一个卷积单元也只处理输入中的局部片段.
        \item 从结构上看, CNN 往往具有三个关键特征:
        \begin{itemize}
            \item \textbf{局部连接}: 每个神经元只接收局部输入, 从而保留空间邻接关系.
            \item \textbf{权重共享}: 相同卷积核在不同空间位置重复使用, 从而降低参数量并增强位置迁移能力.
            \item \textbf{下采样}: 常通过池化或步长大于 $1$ 的卷积减少特征图分辨率, 降低计算开销, 并提升对小范围扰动的稳定性.
        \end{itemize}
        \item 这里还容易把“卷积核大小”和“感受野大小”混为一谈. 对单层卷积而言, 卷积核大小直接决定该层局部连接范围; 但在多层网络中, 高层神经元的有效感受野会随着层数叠加不断扩大, 因而感受野并不等同于某一个卷积核的尺寸.
        \item 从整体上看, CNN 仍然属于前馈神经网络, 只是把上一讲中“线性变换 + 非线性激活”的思想改造成更适合处理空间结构数据的层级表示模型. 后续几节将分别讨论卷积的数学定义、卷积层与池化层的具体形式, 以及典型卷积网络结构.
    \end{itemize}
\end{enumerate}

\textbf{问题思考}:
\begin{itemize}
    \item 若把图像直接展平后送入全连接层, 哪些原本重要的空间结构信息最容易丢失?
    \item 为什么说卷积层更准确地体现的是平移等变性, 而不是单独就实现了完全的平移不变性?
    \item 在多层卷积网络中, 卷积核大小、参数量和感受野增长之间分别存在什么关系?
\end{itemize}

\section{卷积的定义}

卷积最早广泛出现在信号处理和系统分析中, 用来描述“当前输出由当前输入与过去若干时刻输入共同决定”的线性累积过程. 在卷积神经网络中, 卷积被进一步解释为一种局部特征提取操作, 但它的基本数学结构与经典信号处理是一致的.

\begin{enumerate}
    \item \textbf{从延迟叠加理解一维卷积}
    \begin{itemize}
        \item 设离散输入信号为 $\{x_t\}$, 长度为 $K$ 的滤波器系数为 $w_1, \dots, w_K$. 若系统当前输出由最近 $K$ 个时刻的输入加权叠加而成, 则可写为
        \begin{equation*}
            y_t = \sum_{k=1}^{K} w_k x_{t-k+1}.
        \end{equation*}
        这表明: $y_t$ 不是只依赖当前时刻 $x_t$, 而是综合了一个局部时间窗口内的信息.
        \item 例如当 $w_1 = 1, w_2 = \frac{1}{2}, w_3 = \frac{1}{4}$ 时, 有
        \begin{equation*}
            y_t = x_t + \frac{1}{2}x_{t-1} + \frac{1}{4}x_{t-2}.
        \end{equation*}
        这可以理解为: 当前输入保留得最完整, 更早的输入会随着时间衰减.
        \item 无线通信中的多径传播就是一个直观例子. 接收器在时刻 $t$ 除了收到当前路径上的主信号, 还可能收到经过折射、反射或散射后延迟到达的旧信号. 因而输出天然表现为若干延迟分量的叠加, 卷积正好适合刻画这种现象.
        \item 从更一般的角度看, 卷积是线性时不变系统的重要描述工具. 这里的“时不变”意味着: 若输入整体向后平移, 系统输出也会做相同的平移, 这与上一节讨论的平移等变思想是一致的.
    \end{itemize}
    \item \textbf{卷积的数学定义与翻转含义}
    \begin{itemize}
        \item 对一维离散序列 $\boldsymbol{x} \in \mathbb{R}^{N}$ 和卷积核 $\boldsymbol{w} \in \mathbb{R}^{K}$, 常见的离散卷积定义为
        \begin{equation*}
            y_t = {(\boldsymbol{w} * \boldsymbol{x})}_t = \sum_{k=1}^{K} w_k x_{t-k+1}.
        \end{equation*}
        若不做填充且步长为 $1$, 则输出长度为
        \begin{equation*}
            N_{\mathrm{out}} = N - K + 1.
        \end{equation*}
        这也是所谓的窄卷积 (valid convolution).
        \item 很多同学会疑惑: 为什么卷积看起来像是把卷积核“倒过来”再滑动? 从公式上看, 这是因为卷积要比较当前输出位置与输入中相对位移后的各个分量, 因而指标里会出现 $t-k+1$ 这样的反向偏移. 这种写法与线性时不变系统的冲激响应表示天然一致.
        \item 若不对卷积核翻转, 写成
        \begin{equation*}
            y_t = \sum_{k=1}^{K} w_k x_{t+k-1},
        \end{equation*}
        则对应的更准确名称是互相关 (cross-correlation). 现代深度学习框架大多直接实现互相关, 但由于卷积核参数是通过训练学习出来的, 翻转与否通常不会限制模型表达能力, 所以工程上仍习惯称为“卷积层”.
    \end{itemize}
    \item \textbf{步长、零填充与输出尺寸}
    \begin{itemize}
        \item 在实际应用中, 卷积并不一定每次只移动一个位置. 若步长为 $S$, 零填充宽度为 $P$, 则一维情况下输出长度满足
        \begin{equation*}
            N_{\mathrm{out}} = \left\lfloor \frac{N - K + 2P}{S} \right\rfloor + 1.
        \end{equation*}
        其中 $\lfloor \cdot \rfloor$ 表示向下取整. 当 $S = 1$ 且 $P = 0$ 时, 它就退化为前面的 $N-K+1$.
        \item 零填充 (padding) 的作用主要有两个: 一是避免特征图在每一层卷积后都迅速缩小; 二是让边缘位置的像素也能被充分利用. 因此, 现代卷积网络中常采用“same padding”来尽量保持空间分辨率不变.
        \item 步长 (stride) 则控制卷积核滑动的间隔. 步长越大, 输出尺寸越小, 计算量也越低. 需要注意的是, 步长描述的是\textbf{采样间隔}, 而卷积核大小描述的是\textbf{局部感受范围}; 二者都影响输出, 但含义并不相同.
        \item 对二维输入 $\boldsymbol{X} \in \mathbb{R}^{H \times W}$, 卷积核大小为 $U \times V$, 若高和宽方向的填充分别为 $P_h, P_w$, 步长分别为 $S_h, S_w$, 则输出尺寸为
        \begin{equation*}
            H_{\mathrm{out}} = \left\lfloor \frac{H-U+2P_h}{S_h} \right\rfloor + 1,
            \qquad
            W_{\mathrm{out}} = \left\lfloor \frac{W-V+2P_w}{S_w} \right\rfloor + 1.
        \end{equation*}
        这些公式在设计网络结构时非常重要, 因为它们决定了各层特征图能否顺利衔接.
    \end{itemize}
    \item \textbf{滤波器视角: 低通、高通与差分近似}
    \begin{itemize}
        \item 卷积核也可以看作滤波器. 若卷积核更强调邻域平均, 它倾向于保留缓慢变化的低频成分, 这就是低通滤波. 例如
        \begin{equation*}
            \boldsymbol{w}_{\mathrm{lp}} = \frac{1}{3}[1,1,1]
        \end{equation*}
        会把相邻时刻做局部平均, 从而使信号变得更平滑.
        \item 与之相对, 高通滤波更强调快速变化的部分, 常用于突出边缘、尖峰或局部突变. 因此高通滤波不是“磨光”, 而是更偏向放大高频变化. 例如
        \begin{equation*}
            \boldsymbol{w}_{\mathrm{hp}} = [1,-2,1]
        \end{equation*}
        对应二阶差分, 能显著响应曲率变化较大的位置.
        \item 卷积还可近似导数运算. 例如中央差分常写为
        \begin{equation*}
            x'(t) \approx \frac{x_{t+1} - x_{t-1}}{2},
        \end{equation*}
        对应的卷积核可写为 $\frac{1}{2}[-1,0,1]$ 或符号相反的 $\frac{1}{2}[1,0,-1]$. 这里的符号差异通常来自卷积与互相关的定义方向不同, 但它们都在刻画一阶变化率.
        \item 类似地,
        \begin{equation*}
            x''(t) \approx x_{t+1} - 2x_t + x_{t-1}
        \end{equation*}
        对应二阶差分核 $[1,-2,1]$. 这说明卷积并不只是“求加权和”, 它还能够把局部趋势、斜率和曲率等结构信息编码出来.
    \end{itemize}
    \item \textbf{二维卷积与图像特征提取}
    \begin{itemize}
        \item 二维卷积是一维卷积在图像矩阵上的自然推广. 设输入图像为 $\boldsymbol{X} \in \mathbb{R}^{M \times N}$, 卷积核为 $\boldsymbol{W} \in \mathbb{R}^{U \times V}$, 则可写为
        \begin{equation*}
            y_{ij} = \sum_{u=1}^{U} \sum_{v=1}^{V} w_{uv} x_{i-u+1, j-v+1}.
        \end{equation*}
        这表示输出位置 $(i,j)$ 由输入局部邻域与卷积核逐元素加权求和得到.
        \item 在图像处理中, 低通滤波常用于平滑去噪, 例如高斯滤波会保留整体轮廓而削弱随机噪声. 连续情形下, 二维高斯核可写为
        \begin{equation*}
            G(x,y) = \frac{1}{2 \pi \sigma^2} \exp\left(-\frac{x^2 + y^2}{2\sigma^2}\right).
        \end{equation*}
        将它在离散网格上采样并归一化后, 常得到类似如下的 $3 \times 3$ 滤波器:
        \begin{equation*}
            \frac{1}{16}
            \begin{bmatrix}
                1 & 2 & 1 \\
                2 & 4 & 2 \\
                1 & 2 & 1
            \end{bmatrix}.
        \end{equation*}
        之所以称为“高斯滤波”, 是因为滤波器各位置的权重来自高斯函数的取值: 越靠近中心, 权重越大; 越远离中心, 权重按高斯曲线的方式平滑衰减. 这种加权平均方式既能抑制高频噪声, 又不会像简单均值滤波那样过于生硬. 高通滤波则更适合突出边缘和纹理. 例如拉普拉斯算子
        \begin{equation*}
            \begin{bmatrix}
                0 & 1 & 0 \\
                1 & -4 & 1 \\
                0 & 1 & 0
            \end{bmatrix}
        \end{equation*}
        会对像素变化剧烈的位置产生更强响应, 因而可用于边缘检测.
        \item 传统图像处理中, 这些滤波器往往依赖人工设计; 而在卷积神经网络中, 卷积核参数由数据驱动自动学习. 这意味着模型不必预先固定“哪一种边缘最重要”, 而是可以在训练过程中自己发现对任务最有用的局部模式.
        \item 从这个角度看, 卷积既是一个明确的数学运算, 也是连接经典信号处理与现代深度学习的重要桥梁. 下一节讨论的卷积神经网络, 本质上就是把这种局部滤波思想嵌入到多层可学习模型中.
    \end{itemize}
\end{enumerate}

\textbf{问题思考}:
\begin{itemize}
    \item 为什么数学定义中的卷积要涉及卷积核翻转, 而深度学习实现中却常直接使用互相关?
    \item 零填充和步长都会改变输出尺寸, 但它们对特征提取能力的影响有何不同?
    \item 为什么低通滤波更适合平滑噪声, 而高通滤波更容易突出边缘与突变位置?
\end{itemize}

\section{卷积神经网络}

在上一节明确了卷积的数学定义之后, 现在把它放到完整的多层模型中. 卷积神经网络并不是简单地把卷积重复很多次, 而是通过多通道卷积、非线性激活、下采样以及分类头的组合, 逐层把局部模式组织成对任务有用的层级特征.

\begin{enumerate}
    \item \textbf{卷积层的输入输出与互相关实现}
    \begin{itemize}
        \item 设某一卷积层的输入为 $\boldsymbol{X} \in \mathbb{R}^{M \times N \times D}$, 其中 $M, N$ 为空间尺寸, $D$ 为输入通道数. 若该层使用 $P$ 个卷积核, 则卷积核张量可记为 $\boldsymbol{W} \in \mathbb{R}^{U \times V \times D \times P}$, 输出预激活张量为 $\boldsymbol{Z} \in \mathbb{R}^{M' \times N' \times P}$.
        \item 对第 $p$ 个输出通道, 前向传播可写为
        \begin{equation*}
            \boldsymbol{Z}^{(p)} = \sum_{d=1}^{D} \boldsymbol{W}^{(p,d)} \star \boldsymbol{X}^{(d)} + b^{(p)},
            \qquad
            \boldsymbol{Y}^{(p)} = f\left(\boldsymbol{Z}^{(p)}\right).
        \end{equation*}
        这里 $\star$ 表示深度学习实现中更常见的互相关运算, $f$ 通常取 ReLU 等非线性激活函数.
        \item 这说明: 多通道输入并不是先被“压平”成一张二维图再卷积, 而是每个输出通道都要综合所有输入通道的局部响应. 因此, 当输入深度为 $D$ 时, 单个卷积核实际是一个 $U \times V \times D$ 的三维权重块, 而不是彼此独立的若干二维矩阵.
        \item 这里容易与上一节的数学卷积定义混淆. 在严格定义下, 卷积涉及核翻转; 但在 CNN 的工程实现里, 前向传播通常直接采用互相关. 真正决定输出语义的是学习到的参数和多通道聚合方式, 而不是是否显式写出翻转操作.
    \end{itemize}

    \item \textbf{多个卷积核如何形成多通道特征图}
    \begin{itemize}
        \item 若同一层使用 $P$ 个卷积核, 就会得到 $P$ 个输出通道. 不同卷积核可以分别偏向边缘、纹理、颜色组合或更复杂的局部模式, 这也是卷积层表达能力增强的主要来源.
        \item 输入深度 $D$ 决定了每个卷积核需要“看”多厚, 输出深度 $P$ 则由卷积核组数决定. 因而常说的“每个通道对应一组卷积核”更准确地说是: 每个\textbf{输出通道}都对应一组覆盖全部输入通道的权重块.
    \end{itemize}

    \begin{center}
        \begin{tikzpicture}[x=1cm, y=1cm, font=\small, line join=round, line cap=round]
            \draw[blk, fill=black!4, rounded corners=1pt] (0.00,1.10) rectangle (2.30,3.40);
            \draw[blk, fill=black!7, rounded corners=1pt] (0.22,1.30) rectangle (2.52,3.60);
            \draw[blk, fill=black!10, rounded corners=1pt] (0.44,1.50) rectangle (2.74,3.80);
            \node[align=center] at (1.37,0.40) {$\boldsymbol{X}$\\$M\times N\times D$};

            \draw[blk, fill=black!4, rounded corners=1pt] (4.10,3.00) rectangle (4.92,3.82);
            \draw[blk, fill=black!7, rounded corners=1pt] (4.24,3.14) rectangle (5.06,3.96);
            \draw[blk, fill=black!10, rounded corners=1pt] (4.38,3.28) rectangle (5.20,4.10);
            \node[align=center] at (4.65,2.35) {$\boldsymbol{W}^{(1)}$\\$U\times V\times D$};

            \draw[blk, fill=black!4, rounded corners=1pt] (7.00,3.20) rectangle (8.85,5.05);
            \node[align=center] at (7.92,2.55) {$\boldsymbol{Y}^{(1)}$\\$M'\times N'$};

            \draw[blu, fill=blue!6, rounded corners=1pt] (4.10,0.55) rectangle (4.92,1.37);
            \draw[blu, fill=blue!9, rounded corners=1pt] (4.24,0.69) rectangle (5.06,1.51);
            \draw[blu, fill=blue!12, rounded corners=1pt] (4.38,0.83) rectangle (5.20,1.65);
            \node[text=blue!70!black, align=center] at (4.65,-0.10) {$\boldsymbol{W}^{(2)}$\\$U\times V\times D$};

            \draw[blu, fill=blue!6, rounded corners=1pt] (7.00,-0.35) rectangle (8.85,1.50);
            \node[text=blue!70!black, align=center] at (7.92,-0.90) {$\boldsymbol{Y}^{(2)}$\\$M'\times N'$};

            \draw[->, line width=0.9pt] (2.95,3.20) -- (3.95,3.55);
            \draw[->, line width=0.9pt] (5.35,3.55) -- (6.70,3.95);
            \node at (3.45,4.10) {$\star$};
            \node at (6.05,4.35) {$+\,b^{(1)}$};

            \draw[->, line width=0.9pt, draw=blue!70!black] (2.95,1.70) -- (3.95,1.10);
            \draw[->, line width=0.9pt, draw=blue!70!black] (5.35,1.10) -- (6.70,0.88);
            \node[text=blue!70!black] at (3.45,0.85) {$\star$};
            \node[text=blue!70!black] at (6.05,0.15) {$+\,b^{(2)}$};

            \node at (7.92,1.90) {$\vdots$};

            \draw[decorate, decoration={brace, amplitude=4pt}] (9.25,5.05) -- (9.25,-0.35)
                node[midway, right=5pt] {$P$};
            \draw[->, line width=0.9pt] (10.10,2.70) -- (11.20,2.70);

            % \draw[blk, fill=black!4, rounded corners=1pt] (11.50,1.20) rectangle (13.20,2.90);
            \draw[blu, fill=blue!6, rounded corners=1pt] (11.75,1.40) rectangle (13.45,3.10);
            \draw[blk, rounded corners=1pt] (12.00,1.60) rectangle (13.70,3.30);
            \node[align=center] at (12.60,0.35) {$\boldsymbol{Y}$\\$M'\times N'\times P$};
        \end{tikzpicture}

        {\small 图示: 多个卷积核生成多个输出通道, 再沿通道维堆叠.}
    \end{center}

    图中强调了一个容易混淆的点: 输出通道数由卷积核组数决定, 而不是由输入通道数自动继承; 但每一个输出通道都会综合全部输入通道的信息.

    \item \textbf{池化层与下采样的作用}
    \begin{itemize}
        \item 卷积层虽然大幅减少了连接方式和参数量, 但当步长较小、输入分辨率较高时, 特征图的空间尺寸仍然可能很大. 因此, 卷积网络通常还会引入池化层 (pooling layer), 也常称汇聚层, 用于进一步压缩空间维度.
        \item 例如对一个 $4 \times 4$ 特征图做 $2 \times 2$、步长为 $2$ 的最大池化, 可以把它划分为四个不重叠的小块, 并从每个小块中选取最大值, 从而得到一个 $2 \times 2$ 的输出. 这个操作保留了“局部区域中是否出现强响应”的信息, 同时降低了计算量.
        \item 最大池化强调最显著的局部激活, 平均池化则更关注区域内的整体平均水平. 二者都主要作用在空间尺寸上, 通常不会像卷积核个数那样直接改变通道数. 这也是很多同学容易混淆的地方: 卷积层主要负责\textbf{生成新特征}, 池化层主要负责\textbf{压缩空间分辨率}.
        \item 现代网络中也常用步长大于 $1$ 的卷积代替显式池化. 因而“全卷积网络”并不意味着网络失去了下采样能力, 而是把下采样的职责更多交给卷积本身.
    \end{itemize}

    \item \textbf{卷积块与分类头的层级组合}
    \begin{itemize}
        \item 一个典型 CNN 往往由若干“卷积 + 激活 + 下采样”模块堆叠而成, 最后接全连接层或全局平均池化, 再利用 Softmax 将类别得分转化为概率分布. 这一点与第 3 讲的线性分类模型是衔接的: 前面的卷积部分负责构造表征, 后面的分类头负责做判别.
        \item 随着层数加深, 神经元的有效感受野会不断扩大, 特征也会从边缘、纹理逐步过渡到部件和更高层语义. 因此, 深层卷积网络不仅是在“重复同一种滤波”, 而是在做层级表征学习.
        \item 这也解释了为什么现代 CNN 往往偏好“小卷积核 + 更深层数”的设计. 例如连续堆叠多个 $3 \times 3$ 卷积, 可以在保持参数相对可控的同时扩大有效感受野, 并插入更多非线性变换.
        \item LeNet-5 是经典的早期例子: 它通过卷积层、池化层和全连接层的组合实现了手写数字识别, 说明这种分层特征提取思路在实际任务中是有效的. 下一节将进一步讨论这些代表性结构.
    \end{itemize}

    \item \textbf{扩大感受野的方式与空洞卷积}
    \begin{itemize}
        \item 想要让单个神经元“看到”更大的输入区域, 常见方法包括: 增大卷积核尺寸, 增加网络深度, 或者引入空洞卷积 (dilated convolution). 这些方法都会影响感受野, 但代价和适用场景并不相同.
        \item 若空洞率为 $r$, 则一个 $U \times V$ 卷积核对应的有效尺寸会变为
        \begin{equation*}
            U_{\mathrm{eff}} = U + (U - 1)(r - 1),
            \qquad
            V_{\mathrm{eff}} = V + (V - 1)(r - 1).
        \end{equation*}
        这意味着在参数量不变的情况下, 卷积核可以跨过中间若干位置去感受更大范围的上下文.
        \item 需要特别纠正一个常见误解: 空洞卷积的主要作用是\textbf{增大感受野}, 而不是自动“更快地缩小特征图”. 当步长仍为 $1$ 时, 它完全可以保持空间分辨率不变; 若希望尺寸明显缩小, 仍然需要依赖步长或池化.
        \item 因此, 空洞卷积更适合那些既需要大感受野、又不希望过早损失分辨率的任务, 例如语义分割或密集预测.
    \end{itemize}

    \item \textbf{转置卷积与上采样需求}
    \begin{itemize}
        \item 在分类任务中, 网络通常随着层数加深而逐步缩小空间尺寸; 但在图像生成、超分辨率和语义分割等任务中, 我们常常又需要把低分辨率特征恢复为高分辨率输出. 这时就会用到转置卷积, 也常被旧教材称为微步卷积.
        \item 从线性代数角度看, 转置卷积对应的是卷积算子的转置, 它并不是普通卷积的真正逆运算. 这一点很重要: 输入经过一次卷积后丢失的信息, 并不会仅靠“转置”就被自动恢复回来.
        \item 从实现角度看, 可以把转置卷积理解为: 先在输入元素之间插入空位或零值, 再做一次普通卷积, 从而增大输出尺寸. 这与空洞卷积在形式上有相似之处, 但两者的目的相反: 空洞卷积是在\textbf{不降采样的情况下扩大感受野}, 转置卷积是在\textbf{上采样时重建空间网格}.
    \end{itemize}
\end{enumerate}

\subsection{卷积层的反向传播}

卷积层的训练仍然遵循第 4 讲中的链式法则, 只是线性映射从矩阵乘法变成了带权共享的局部运算. 为了突出核心结构, 下面先以单通道、valid 互相关为例:
\begin{equation*}
    y_{ij} = \sum_{u=1}^{U} \sum_{v=1}^{V} W_{uv} X_{i+u-1,\, j+v-1},
    \qquad
    L = f(Y).
\end{equation*}

\begin{enumerate}
    \item \textbf{卷积核的梯度}
    \begin{itemize}
        \item 由链式法则可得
        \begin{equation*}
            \frac{\partial L}{\partial W_{uv}}
            =
            \sum_{i,j}
            \frac{\partial L}{\partial y_{ij}}
            X_{i+u-1,\, j+v-1}
            =
            {\left(\frac{\partial L}{\partial Y} \star X\right)}_{uv}.
        \end{equation*}
        \item 这表明卷积核梯度本质上是“输出误差图”与输入特征之间的互相关. 直观地说, 若某个局部输入模式经常与较大的误差同时出现, 那么相应权重就会被重点更新.
    \end{itemize}

    \item \textbf{输入的梯度}
    \begin{itemize}
        \item 对输入求导可得
        \begin{equation*}
            \frac{\partial L}{\partial X}
            =
            \frac{\partial L}{\partial Y} * \tilde{W},
            \qquad
            \tilde{W}_{uv} = W_{U-u+1,\, V-v+1}.
        \end{equation*}
        \item 这里出现了卷积核翻转, 本质上来自前向线性算子的转置. 因而“前向传播常用互相关, 反向传播为什么又会出现翻转”并不矛盾; 它只是说明输入梯度要把误差重新分配回原来的空间位置.
    \end{itemize}

    \item \textbf{多通道情形}
    \begin{itemize}
        \item 当输入为 $\boldsymbol{X} \in \mathbb{R}^{M \times N \times D}$, 输出预激活为 $\boldsymbol{Z} \in \mathbb{R}^{M' \times N' \times P}$ 时, 第 $p$ 个输出通道和第 $d$ 个输入通道之间的梯度关系可写为
        \begin{align*}
            \frac{\partial L}{\partial \boldsymbol{W}^{(p,d)}} &= \frac{\partial L}{\partial \boldsymbol{Z}^{(p)}} \star \boldsymbol{X}^{(d)}, \\
            \frac{\partial L}{\partial b^{(p)}} &= \sum_{i,j} \frac{\partial L}{\partial Z_{ij}^{(p)}}, \\
            \frac{\partial L}{\partial \boldsymbol{X}^{(d)}} &= \sum_{p=1}^{P} \frac{\partial L}{\partial \boldsymbol{Z}^{(p)}} * \tilde{\boldsymbol{W}}^{(p,d)}.
        \end{align*}
        \item 最后一式说明: 同一个输入通道往往会被多个输出通道共同使用, 因而其梯度需要把来自所有输出通道的贡献加总起来. 这正是权重共享和多通道耦合在反向传播中的体现.
    \end{itemize}
\end{enumerate}

\textbf{问题思考}:
\begin{itemize}
    \item 为什么说单个输出通道对应的是一个覆盖全部输入通道的三维卷积核, 而不是若干彼此独立的二维核?
    \item 池化、步长卷积和空洞卷积都会影响感受野或输出尺寸, 它们三者分别解决的核心问题是什么?
    \item 在卷积层反向传播中, 为什么对输入求梯度时会出现卷积核翻转, 而对卷积核求梯度时更像是在做互相关?
\end{itemize}
\section{卷积神经网络的例子}
\begin{enumerate}
    \item \textbf{LeNet-5: 卷积网络的早期范式}
    \begin{itemize}
        \item LeCun 等人在 1989--1998 年间逐步发展出 LeNet 系列, 其中 LeNet-5 是最经典的版本, 主要用于邮政编码和银行支票中的手写数字识别. 它说明了卷积、下采样与全连接分类头可以组成一个端到端可训练的视觉模型.
        \item LeNet-5 的主干结构可概括为“卷积-下采样-卷积-下采样-全连接-输出”. 例如对 $32 \times 32$ 输入做一次 $5 \times 5$ 的 valid 卷积后, 空间尺寸会先变成 $28 \times 28$, 再经过下采样缩小. 这正好对应上一节讨论的输出尺寸公式.
        \item 从今天的角度看, LeNet-5 的规模并不大, 但它首次系统地展示了局部连接和权重共享在视觉任务中的优势, 因而常被视为现代卷积神经网络的起点.
    \end{itemize}

    \item \textbf{AlexNet: 深度学习复兴的标志}
    \begin{itemize}
        \item 20 世纪 90 年代后, 视觉领域一度更多依赖人工特征与浅层分类器. 直到 2012 年, Alex Krizhevsky, Ilya Sutskever 和 Geoffrey Hinton 提出的 AlexNet 在 ILSVRC 中把 top-5 错误率从上一代方法的大约 $26.2\%$ 降到 $15.3\%$, 卷积网络才重新成为主流.
        \item AlexNet 并不是第一个 CNN, 但它是第一个在大规模数据集和现代硬件条件下取得压倒性优势的深层卷积网络. 它使用了 ReLU 激活、数据增强、Dropout 和 GPU 并行训练等关键技术, 这些做法后来几乎成为视觉模型训练的标准配置.
        \item 其网络结构大致包含 5 个卷积层和 3 个全连接层. 当时由于单张 GPU 显存有限, AlexNet 采用双 GPU 分路训练. 这里要注意, 这不是“模型必须拆成两半”的理论要求, 而是受当时硬件条件约束的工程实现.
    \end{itemize}

    \item \textbf{Inception: 多尺度并行特征提取}
    \begin{itemize}
        \item GoogLeNet / Inception 在 2014 年的 ILSVRC 中取得优异结果. 它的核心思想不是一味加深或加宽网络, 而是在同一层内并行放置 $1 \times 1, 3 \times 3, 5 \times 5$ 卷积和池化分支, 再把结果沿通道维拼接起来.
        \item 这样做的好处在于, 不同分支可以同时捕捉不同尺度的局部模式. 小卷积核更擅长细粒度纹理, 大卷积核更擅长较大范围的结构信息, 池化分支则提供一定的局部不变性.
        \item 这里容易误解成“把不同卷积核简单堆叠起来”. 更准确地说, Inception 模块采用的是\textbf{并行分支 + 通道拼接}. 其中 $1 \times 1$ 卷积还常被用作瓶颈层, 先做通道混合与降维, 再减少后续大卷积核的计算量.
        \item 后续版本如 Inception-v3 还把 $N \times N$ 卷积分解为 $N \times 1$ 与 $1 \times N$, 进一步降低参数和计算成本.
    \end{itemize}

    \item \textbf{ResNet: 用残差连接训练更深网络}
    \begin{itemize}
        \item 2015 年提出的 ResNet 说明, 网络变深后真正困难的不只是参数增多, 还有优化难度显著上升. 一个典型残差块可写为
        \begin{equation*}
            \boldsymbol{y} = \mathcal{F}(\boldsymbol{x}) + \boldsymbol{x},
        \end{equation*}
        其中 $\mathcal{F}(\boldsymbol{x})$ 表示若干卷积层学习到的残差映射.
        \item 这意味着网络不必在每个模块里都重新学习完整目标映射, 而可以转而学习“相对输入还需要补充什么”. 若最优映射接近恒等映射, 那么让 $\mathcal{F}(\boldsymbol{x})$ 接近 0 往往比直接拟合一个复杂变换更容易.
        \item 残差连接确实有助于缓解梯度消失, 因为误差信号可以沿着捷径路径更直接地传回前层; 但它更核心的作用是缓解深层网络的退化问题, 让 50 层、101 层甚至 152 层网络都能稳定优化. 因而不能把 ResNet 简单理解成“只为了解决梯度太小”.
    \end{itemize}
\end{enumerate}

\textbf{问题思考}:
\begin{itemize}
    \item 为什么说 AlexNet 是卷积网络复兴的标志, 但并不是第一个卷积神经网络?
    \item Inception 中并行分支和 ResNet 中残差连接, 分别在解决哪一类结构设计问题?
    \item 若把 LeNet-5, AlexNet 和 ResNet 放在同一条发展线上看, 它们分别说明了卷积网络的哪些关键能力?
\end{itemize}

\section{卷积神经网络的应用}
\begin{enumerate}
    \item \textbf{AlphaGo 与棋盘决策中的空间表征}
    \begin{itemize}
        \item 卷积网络并不只适用于自然图像. 在围棋中, 棋盘同样具有规则的二维拓扑结构, 因而可以把局面表示为若干特征平面的堆叠. AlphaGo 正是利用这一点, 让网络从 $19 \times 19$ 棋盘中抽取局部落子形状、联络关系和全局攻防格局.
        \item AlphaGo 在 2015 年击败欧洲冠军 Fan Hui, 并在 2016 年与李世石的人机对局中广为人知. 其核心并不是“单独一个 CNN”, 而是策略网络、价值网络与蒙特卡罗树搜索共同组成的系统: 策略网络给出下一步落子的概率分布, 价值网络估计当前局面的胜率, 搜索过程再利用这些信息筛选更优分支.
        \item 从本课程角度看, AlphaGo 的意义在于说明: 只要数据存在明确的局部空间结构, 卷积网络就可以作为强有力的表征学习模块, 即使最终任务并不是传统图像分类.
    \end{itemize}

    \item \textbf{目标检测、实例分割与 OCR}
    \begin{itemize}
        \item 在目标检测中, 模型不仅要判断“图中有什么”, 还要回答“它在哪里”. 因此输出通常同时包含类别标签和边界框位置. YOLO, Faster R-CNN 等模型都建立在卷积特征图之上, 再接检测头完成定位与分类.
        \item 实例分割比目标检测更进一步, 它要求对每个目标给出像素级掩码. 例如 Mask R-CNN 可以看作“检测 + 分割”联合建模: 先找出候选目标, 再在局部区域内预测精细轮廓.
        \item OCR 则把卷积网络用于文字图像理解. 卷积层先提取笔画、边缘和局部字形结构, 再交给序列解码或分类模块输出字符. 因而 OCR 不是简单的“图片分类”, 而是局部视觉模式识别与语言约束结合的结果.
    \end{itemize}

    \item \textbf{特征可视化、风格迁移与图像生成}
    \begin{itemize}
        \item Deep Dream 利用预训练卷积网络的中间激活来“反向放大”某些特征. 具体地, 它固定网络参数, 反过来优化输入图像, 使选定层的激活更强. 这样生成的图像常呈现出夸张的纹理和重复图案, 有助于理解 CNN 在不同层学到了什么.
        \item 神经风格迁移则把卷积特征进一步用于“内容”和“风格”的分离表示. 内容部分强调高层语义结构, 风格部分常通过特征相关性或 Gram 矩阵刻画纹理统计. 于是模型可以在保留内容轮廓的同时, 把另一幅图的绘画风格迁移过来.
        \item 现代“文生图”系统已经超出了早期纯 CNN 生成模型的范畴. 一个典型流程是: 先由文本编码器把提示词映射为语义向量, 再由扩散模型在潜空间中逐步去噪生成图像表示, 最后用解码器还原为像素图像. 其中 U-Net 式去噪器仍然保留了卷积网络擅长建模局部结构的思想, 但整体系统通常还会结合注意力机制来处理长距离语义对齐.
    \end{itemize}

    \item \textbf{对抗样本与鲁棒性分析}
    \begin{itemize}
        \item 卷积网络在实际部署中并非总是可靠. 有时只需对输入加入极小扰动, 人眼几乎察觉不到变化, 模型的预测却可能发生明显偏移. 这说明高维输入空间中的决策边界可能比直觉更脆弱.
        \item 以 FGSM 为例, 对抗样本可写为
        \begin{equation*}
            \boldsymbol{x}_{\mathrm{adv}} = \boldsymbol{x} + \epsilon \cdot \operatorname{sign}\!\left(\nabla_{\boldsymbol{x}} \mathcal{L}(\theta, \boldsymbol{x}, y)\right).
        \end{equation*}
        这里 $\epsilon$ 控制扰动幅度, $\nabla_{\boldsymbol{x}} \mathcal{L}$ 表示损失函数对输入的梯度. 该方法的直观含义是: 在输入空间中沿着“最能增大损失”的方向迈出一小步, 就可能把样本推过分类边界.
        \item 因而对抗样本研究并不是为了证明 CNN “完全没用”, 而是为了分析其脆弱性并改进鲁棒性. 常见改进思路包括对抗训练、输入去噪、结构正则化以及更稳定的训练目标设计.
    \end{itemize}
\end{enumerate}

需要说明的是, 卷积同样可以用于文本中的局部 $n$-gram 模式提取, 但由于后续课程将专门讨论 Transformer 与序列建模, 这里不再展开.

\textbf{问题思考}:
\begin{itemize}
    \item 为什么 AlphaGo 这样的系统虽然核心任务是决策, 仍然需要卷积网络来处理棋盘状态?
    \item 目标检测、实例分割和 OCR 都依赖卷积特征, 但它们在输出形式上有何本质区别?
    \item 现代文生图系统与早期基于 CNN 的 Deep Dream 或风格迁移相比, 在生成机制上发生了哪些关键变化?
\end{itemize}

注: 详细定义和更多内容请参考课程教材第 5 章.

\end{document}
